\documentclass[11pt]{article}
\usepackage[margin=1in]{geometry}
\usepackage{amsmath,amssymb,amsthm}
\usepackage{hyperref}
\usepackage{enumitem}

\newtheorem{theorem}{Theorem}
\newtheorem{lemma}[theorem]{Lemma}
\newtheorem{corollary}[theorem]{Corollary}
\theoremstyle{definition}
\newtheorem{definition}[theorem]{Definition}
\theoremstyle{remark}
\newtheorem{remark}[theorem]{Remark}

\usepackage{xurl}
\let\oldus\_
\renewcommand{\_}{\oldus\allowbreak}
\newcommand{\Fix}{\operatorname{Fix}}
\newcommand{\per}{\operatorname{per}}

\title{A disproof of Conjecture 00000005400\\[2pt]
\large Conjugate maps cannot have different periods}
\author{Jackmeson1}
\date{}

\begin{document}
\sloppy
\maketitle

\section{The conjecture}

Conjecture 00000005400 reads (English version; the Chinese version says the same, with
``explicit conjugate pair with the same measure and different periods'' for the last clause):

\begin{quote}
\emph{Definition:} Orbit distribution and period counting are two layers.
\emph{Conjecture:} There exist two maps with identical orbit distribution limits but different
periodic point counts, and the separation is realized by an explicit conjugate pair with the same
measure but different periods.
\end{quote}

It is a conjunction. Besides two maps $f,g$ with (A) identical orbit distribution limits and
(B) different periodic point counts, it requires (C) an explicit \emph{conjugate pair} with the same
measure and \emph{different periods}. We show that (C) can never hold, whatever (A) and ``the same
measure'' mean. Hence the conjunction is false.

\section{Reading}

\begin{definition}
Maps $f\colon X\to X$ and $g\colon Y\to Y$ are \emph{conjugate} if there is a bijection
$h\colon X\to Y$ with $h\circ f=g\circ h$, i.e.\ $h(f(x))=g(h(x))$ for every $x\in X$. A
\emph{topological} conjugacy is such an $h$ that is a homeomorphism; a \emph{measure-preserving}
conjugacy is such an $h$ that is a measurable isomorphism with $h_*\mu=\nu$.
\end{definition}

This is the standard meaning of ``conjugate maps'' in dynamics (e.g.\ Katok--Hasselblatt,
\emph{Introduction to the Modern Theory of Dynamical Systems}), and it is the reading under which the
sibling conjecture 00000005390 (``\ldots realized by an explicit conjugate pair \ldots'') was refuted in
an accepted submission. Every topological or measure-preserving conjugacy is in particular a bijective
conjugacy, so refuting the bijective case refutes all of them. The conjecture names ``the same measure''
as a separate condition on the pair; we keep it, as an arbitrary predicate on $(f,g,h)$. An isomorphism
of measure-preserving systems defined only almost everywhere is a different notion: it ignores null
sets and therefore says nothing about individual periodic points, whereas the conjecture asks the
conjugate pair to differ precisely in \emph{periods}, the point-level ``period counting'' layer.

\emph{Periodic data.} We use the standard notions (as in Mathlib): $x$ is $n$-periodic if
$f^n(x)=x$; $\Fix(f^n)$ is the set of such $x$; the \emph{least (minimal) period} $\per_f(x)$ is the
least $n\ge1$ with $f^n(x)=x$; $P_n(f)=\{x \text{ periodic}: \per_f(x)=n\}$; an $n$-cycle is the periodic
orbit (as a cyclic list) of a point of $P_n(f)$; and the period set is $\{\per_f(x): x\text{ periodic}\}$.
Counts are cardinalities, so infinite sets are allowed. ``Different periodic point counts / different
periods'' is read in the broadest way: \emph{some} $|\Fix(f^n)|$, $|P_n(f)|$ or number of $n$-cycles
differs, or the period sets differ, or (for the conjugate pair) some point $x$ has
$\per_g(h(x))\neq\per_f(x)$. Any narrower reading is a special case. We also allow the conjugate pair
$(f',g')$ to be different from $(f,g)$; the reading in which $(f,g)$ itself is the conjugate pair is a
special case (Corollary~\ref{cor:same}).

\section{Conjugacy preserves all periodic data}

Fix a conjugacy $h\colon X\to Y$, $h\circ f=g\circ h$, with $h$ bijective.
\emph{Convention.} $\per_f(x)$ denotes the least positive $n$ with $f^n(x)=x$ if such an $n$ exists, and
$\per_f(x)=0$ if $x$ is not periodic; this is exactly Mathlib's \texttt{Function.minimalPeriod}. With
this convention the equality of least periods below holds at every point, periodic or not.

\begin{lemma}\label{lem:iter}
For every $n\ge0$ and $x\in X$: $h(f^n(x))=g^n(h(x))$, hence $f^n(x)=x \iff g^n(h(x))=h(x)$.
\end{lemma}
\begin{proof}
Induction on $n$ using $h\circ f=g\circ h$; the equivalence follows from injectivity of $h$, since
$g^n(h(x))=h(x)$ means $h(f^n(x))=h(x)$.
\end{proof}

\begin{lemma}\label{lem:per}
$x$ is periodic for $f$ iff $h(x)$ is periodic for $g$, and $\per_g(h(x))=\per_f(x)$ for every $x$.
\end{lemma}
\begin{proof}
By Lemma~\ref{lem:iter} the sets $\{n: f^n(x)=x\}$ and $\{n: g^n(h(x))=h(x)\}$ coincide, so their
least positive elements coincide (or both do not exist).
\end{proof}

\begin{lemma}\label{lem:bij}
For every $n$, $h$ restricts to bijections $\Fix(f^n)\to\Fix(g^n)$ and $P_n(f)\to P_n(g)$; the map
$c\mapsto h(c)$ is a bijection from the $n$-cycles of $f$ to the $n$-cycles of $g$; and the period
sets of $f$ and $g$ are equal.
\end{lemma}
\begin{proof}
The first two statements follow from Lemmas~\ref{lem:iter} and~\ref{lem:per}, since $h$ is a bijection
of $X$ onto $Y$ mapping each set onto the other. For cycles: the orbit of $x$ is the cyclic list
$(x,f(x),\dots,f^{p-1}(x))$ with $p=\per_f(x)$; applying $h$ termwise gives
$(h(x),g(h(x)),\dots,g^{p-1}(h(x)))$, which is the orbit of $h(x)$ because $\per_g(h(x))=p$. So $h$ maps
$n$-cycles to $n$-cycles. Since $h^{-1}$ conjugates $g$ to $f$, the same holds for $h^{-1}$, and the two
maps are mutually inverse on cycles. The period sets agree by Lemma~\ref{lem:per} and surjectivity of $h$.
\end{proof}

\begin{theorem}\label{thm:main}
For every interpretation of ``identical orbit distribution limits'' and ``the same measure'',
Conjecture 00000005400 is false.
\end{theorem}
\begin{proof}
Clause (C) asks for a conjugate pair $(f',g')$, with conjugacy $h$, such that some periodic count of
$f'$ and $g'$ differs, or the period sets differ, or $\per_{g'}(h(x))\ne\per_{f'}(x)$ for some $x$.
Lemma~\ref{lem:bij} gives equal cardinalities (via explicit bijections) and equal period sets, and
Lemma~\ref{lem:per} gives equal least periods at every point. So (C) fails, whatever extra condition is
imposed on $(f',g',h)$, and so does the conjunction.
\end{proof}

\begin{corollary}\label{cor:same}
There are no maps $f,g$ and conjugacy $h$ such that $(f,g)$ satisfies (A), (B), the same-measure
condition, and $h\circ f=g\circ h$. The same holds when $h$ is a homeomorphism, or a measure-preserving
measurable isomorphism ($h_*\mu=\nu$, conjugacy at every point).
\end{corollary}

Clauses (A) and (B) on their own are not disputed here: two \emph{non-conjugate} maps may well share
orbit-distribution limits and differ in periodic counts. What fails is the conjecture's requirement
that the separation be realized by a conjugate pair with different periods.

\section{Lean 4 formalization}

The file \texttt{lean/Conjecture5400/Basic.lean} (Lean 4.33.1, Mathlib v4.33.1) uses Mathlib's
\texttt{Function.Semiconj h f g} ($h\circ f=g\circ h$ pointwise) with \texttt{h : X $\simeq$ Y},
\texttt{Function.IsPeriodicPt}, \texttt{Function.ptsOfPeriod} ($\Fix(f^n)$),
\texttt{Function.periodicPts}, \texttt{Function.minimalPeriod} ($\per_f$) and
\texttt{Function.periodicOrbit} (cycles as Mathlib \texttt{Cycle}s). Counts are \texttt{Cardinal.mk}.

\begin{itemize}[leftmargin=*]
\item \texttt{isPeriodicPt\_iff}, \texttt{minimalPeriod\_eq}: Lemmas~\ref{lem:iter} and~\ref{lem:per}.
\item \texttt{ptsOfPeriodEquiv}, \texttt{leastPeriodPtsEquiv}, \texttt{cyclesEquiv},
\texttt{periodSet\_eq}: Lemma~\ref{lem:bij}, as explicit \texttt{Equiv}s.
\item \texttt{not\_periodDataDiffer}: a conjugate pair has no difference in periodic data and equal
least periods at corresponding points.
\item \texttt{Statement SameOrbitLimits SameMeasure}: the conjecture, with the two side conditions as
arbitrary predicates (\texttt{SameMeasure} may depend on $h$).
\item \texttt{conjecture5400\_false}: $\neg$\texttt{Statement SameOrbitLimits SameMeasure} for all
such predicates (Theorem~\ref{thm:main}).
\item \texttt{conjecture5400\_samePair\_false}, \texttt{homeomorph\_conjugate\_same\_periods},
\texttt{measured\_conjugate\_same\_periods}: Corollary~\ref{cor:same}.
\end{itemize}

There is no \texttt{sorry}, no \texttt{native\_decide} and no added axiom;
\texttt{\#print axioms Conjecture5400.conjecture5400\_false} reports only \texttt{propext},
\texttt{Classical.choice} and \texttt{Quot.sound}. The file also elaborates with
\texttt{-DwarningAsError=true}. To reproduce, in \texttt{lean/}:
\begin{verbatim}
lake exe cache get
lake build
\end{verbatim}

\section{Credits and prior work}

The transport argument (iterates intertwined, periodic and exact-period sets in bijection, cycles
mapped by $h$, arbitrary extra property on the conjugacy) is adapted from the accepted disproof of
the sibling conjecture 00000005390 by ziangni-sys
(\url{solutions/00000005390/ziangni-sys_submission_20261004205500}, The Last Math Competition,
GPL-3.0). Our Lean file is rewritten on top of Mathlib's periodic-point API and states the conjecture
00000005400 itself.

An earlier submission for 00000005400 (PR \#25) was removed from the repository on 2026-10-01 as having
``refuted a mechanism/reading, not the literal statement''; its Lean enumerated permutations of a
4-point set. The present disproof is not restricted to finite sets or permutations: it covers
all maps on all spaces, all cardinalities, every periodic count, least periods and cycles, and leaves
the side conditions of the literal statement completely general.

\end{document}
